% TOP_PROBLEM: {"id": 2853, "title": "Kirby Problem 3.55", "category_id": 7, "category": {"id": 7, "name": "topology", "display_name": "Topology", "description": "Properties preserved under continuous deformations.", "slug": "topology", "order_index": 7, "created_at": "2026-07-31T15:26:25.671Z"}, "status": "open"}
% FIRST_POSTED: null
% ENTRY_KIND: statement_only
% Imported from: batch12.tex; UnsolvedMath ID 2853
\documentclass[11pt]{article}
\usepackage[margin=25mm]{geometry}
\usepackage{fontspec}
\setmainfont{Latin Modern Roman}
\usepackage{amsmath,amssymb,mathtools}
\usepackage{unicode-math}
\setmathfont{Latin Modern Math}
\usepackage{xurl,hyperref}
\hypersetup{colorlinks=true,urlcolor=blue}
\setcounter{secnumdepth}{0}
\setlength{\parindent}{0pt}
\setlength{\parskip}{5pt}
\setlength{\emergencystretch}{3em}
\newcommand{\sref}[2]{\hyperlink{source:#1:#2}{[S#2]}}
\newcommand{\cataloguescope}[1]{\paragraph{Recorded target.}#1}
\begin{document}
% UnsolvedMath ID 2853; source catalogue code KP-3.55
\setcounter{section}{2852}
\section{Kirby Problem 3.55}\label{2853}
{\small\sffamily UnsolvedMath ID 2853 · Topology}
\subsection{Definitions and mathematical statement}

\paragraph{Shared notation.}$\mathbb N=\{1,2,\ldots\}$, and $\mathbb Z,\mathbb Q,\mathbb R,\mathbb C$ denote the integers, rationals, reals and complex numbers. Any other conventions are specified locally.


Let $R=\mathbb F_2[U]$, where $\mathbb F_2=\mathbb Z/2\mathbb Z$. For a knot $K\subset S^3$, $\operatorname{HFK}^-(K)$ is minus knot Floer homology over $R$, retaining its $U$ action. For a rational homology $3$-sphere $Y$, define
\[
\operatorname{HF}_{\rm red}(Y)=\operatorname{coker}\bigl(\operatorname{HF}^{\infty}(Y;\mathbb F_2)\longrightarrow\operatorname{HF}^{+}(Y;\mathbb F_2)\bigr),
\]
with its induced $R$ action. All spin$^c$ structures are included. To admit an $\mathbb F_2$ summand as an $R$-module means $M\cong R/(U)\oplus N$ for some $R$-module $N$; a one-dimensional submodule alone is insufficient.
\paragraph{Statement recorded in the supplied source.}
(a) Does every nontrivial knot $K\subset S^3$ have an $R/(U)$ direct summand in $\operatorname{HFK}^-(K)$? (b) Does every rational homology sphere $Y$ with $\operatorname{HF}_{\rm red}(Y)\ne0$ have an $R/(U)$ direct summand in that reduced module?
The source also asks about a geometric interpretation of the summand's generator and records the refinement: if $\operatorname{HF}_{\rm red}(Y)$ has a direct summand $R/(U^k)$, must it have a direct summand $R/(U^\ell)$ for each $1\le\ell<k$? \sref{2853}{2}
\subsection{Short English statement}
Must minus knot Floer homology of every nontrivial knot, and every nonzero reduced Floer group of a rational homology sphere, split off a one-dimensional summand annihilated by U? What geometric meaning and smaller-length summands accompany such splittings?
\subsection{Sources}
\begin{description}
\item[\hypertarget{source:2853:1}{S1}] ULAM.AI and the UnsolvedMath Contributors, \emph{UnsolvedMath: A Curated Collection of Open Mathematics Problems}, record 2853 (KP-3.55). CC BY 4.0. \url{https://huggingface.co/datasets/ulamai/UnsolvedMath}.
\item[\hypertarget{source:2853:2}{S2}] R. İnanç Baykur, Robion C. Kirby and Daniel Ruberman (editors), \emph{K3: A New Problem List in Low-Dimensional Topology}, Mathematical Surveys and Monographs 295 (2026), Problem 3.55, p. 170 and following remarks; original author version. \url{https://math.berkeley.edu/sites/default/files/surv-295-ruberman-watermarked-author-pdf.pdf}.
\end{description}
\label{end:2853}
\end{document}
